% Section 8: Implications

\section{Implications}
\label{sec:implications}

\subsection{For Law Firms}
\label{sec:law-firms}

\paragraph{Document Review Tools:} The cross-reference specification suggests a practical tool: automated flagging of definition inconsistencies across fund documents. Even without full formalization, a system that highlights divergent terms would reduce review time and catch errors.

\paragraph{Template Verification:} Law firms maintain precedent banks of fund documents. Catala-style specifications could verify that template amendments don't break existing logic---that adding a new investor category doesn't create orphan cross-references.

\paragraph{Training:} Specifications could help train junior associates. Rather than reading 100 pages of LPA prose, a trainee could study the logical structure: what are the conditions for GP removal? What triggers an LPAC meeting? The specification is a map of the document's logic.

\subsection{For Regulators}
\label{sec:regulators}

\paragraph{Machine-Readable Regulation:} Could CSSF circulars be published in Catala alongside natural language? This would eliminate interpretation disputes and allow automated compliance checking. The DG TAXUD work on computational tax illustrates the potential.

\paragraph{Partial Automation:} Pre-screening of RAIF filings for structural compliance (authorized AIFM? depositary appointed?) could be automated, freeing reviewers for substantive assessment.

\paragraph{Gap Identification:} Formal specification often reveals gaps in regulation. The encoding process forces clarity: what \textit{exactly} does ``significant size'' mean? This could inform regulatory revision.

\subsection{For the Industry}
\label{sec:industry}

\paragraph{Standardization Pressure:} Verification works best with consistent structures. If the industry adopted standard definitions (perhaps through ILPA or Invest Europe), cross-document consistency checking becomes trivial.

\paragraph{Cost-Benefit:} Full specification requires significant upfront investment. The question is: does the inconsistency detection value exceed the encoding cost? For €400M+ funds with complex side letter stacks, possibly yes.

\paragraph{Liability Questions:} If a fund relies on formal verification and the specification is wrong, who bears liability? The specifier? The law firm that reviewed it? This creates new professional responsibility questions.

\subsection{For Computational Law Research}
\label{sec:research}

This paper demonstrates that:

\begin{enumerate}
  \item Catala's default logic handles the fund law exception hierarchy naturally
  \item Real inconsistencies in practice (like the BBB- vs.\ ``investment grade'' example) are detectable
  \item The limits of formalization are identifiable and documentable
\end{enumerate}

Future research could extend to other regulatory domains: UCITS rules, MiFID product governance, PRIIPs disclosure requirements.
